\documentclass{article}
\usepackage{amsmath,amssymb}
\usepackage{booktabs}
\usepackage[margin=1in]{geometry}
\title{Step 230: Hilbert--Schmidt Trace Diagnostic for \(C_\ell P_\eta\)}
\author{Six Birds Foundations III}
\date{2026-05-16}
\begin{document}
\maketitle

\section{Target}

The tested operator is
\[
  C_\ell P_\eta = (I-P_\infty)M_{m_\ell}P_\infty P_\eta
\]
where \(P_\eta\) projects onto
\[
  H_\eta=\overline{\operatorname{span}}\{\kappa_{\rho_i,k}^{1/2}\}.
\]
With simple zeta zeros this becomes the countable span of the zero-order vectors \(\kappa_{\rho_i,0}^{1/2}\).

\section{Hilbert--Schmidt Diagnostic}

For an orthonormal basis \(\{v_n\}\) of \(H_\eta\),
\[
  \|C_\ell P_\eta\|_{\rm HS}^2
  =
  \sum_n \|C_\ell v_n\|^2.
\]
If this sum is finite, \(C_\ell P_\eta\) is Hilbert--Schmidt and hence compact. If it diverges, the operator is not Hilbert--Schmidt, although compactness is not thereby decided.

\section{Computation}

The computation reuses the Step 216 Gram--Schmidt data generated from the Step 214/215 finite PSWF/sinc model. CAND1 is the Burnol-boundary model; CAND2 is the zeta-form comparison model.

\begin{center}
\begin{tabular}{lrrrr}
\toprule
model & tested basis & finite terms & usable terms & partial sum\\
\midrule
CAND1 & 20 & 15 & 2 & 8.8864712371\\
CAND2 & 10 & 10 & 10 & 2.9957793326\\
\bottomrule
\end{tabular}
\end{center}

The CAND1 tail contains apparent nondecaying positive terms, but they occur after the Gram--Schmidt residual norms have collapsed to low absolute signal. Hence they cannot certify divergence. CAND2 remains stable and nondecaying, but it is not the verified \(a=1/2\) transport formula.

\section{Verdict}

\[
  \boxed{\texttt{V\_HS\_partial}.}
\]

The HS route remains open. The next required input is a stable infinite orthonormal basis model for \(H_\eta\) under the verified Burnol transport, or a theorem replacing finite-grid Gram--Schmidt.

\end{document}
